\Needspace{8\baselineskip}
\Article{258}
All the forms found in the preceding article clearly belong to ambiguous classes, and conversely every properly primitive ambiguous class of determinant \(D\) must contain at least one of those forms. Such a class certainly contains ambiguous forms, and every properly primitive ambiguous form \((a,b,c)\) of determinant \(D\) is equivalent to one of the forms of the preceding article, namely either
\[
(a,0,-\tfrac{D}{a})\text{ or }(a,\tfrac12 a,\tfrac14 a-\tfrac{D}{a})
\]
according as \(b\) is \(\equiv0\) or \(\equiv\frac12 a\ (\mathrm{mod.}\,a)\). The problem is therefore reduced to finding how many distinct classes those forms constitute.

If the form \((a,0,c)\) is among the forms of the preceding article, the form \((c,0,a)\) occurs among them too and is always different from it, with the single exception \(a=c=\pm1\), and therefore \(D=-1\), which we will set aside for a while. Since these forms clearly belong to the same class, it is enough to retain one, and we will discard the one whose leading coefficient is greater than the third; we also set aside the case \(a=-c=\pm1\), or \(D=1\). In this way we can reduce all the forms
\SourceLeaf{292}
$(A,0,C)$ to half their number, always retaining one of each pair; and in all those remaining we have $A<\sqrt{\pm D}$.

Likewise, if the form $(2b,b,c)$ occurs among the forms of the preceding article, so does
\[
(4c-2b,\ 2c-b,\ c)=\left(-\frac{2D}{b},\ -\frac{D}{b},\ c\right)
\]
which is properly equivalent to it and different from it, with the single exception that we set aside, $c=b=\pm1$, or $D=-1$. Of these two forms it is enough to retain the one whose leading coefficient is less than the other's (in this case they cannot be equal in magnitude but different in sign). Thus all the forms $(2B,B,C)$ too can be reduced to half their number, always discarding one of each pair; and in those remaining we have $B<\frac{D}{B}$, or $B<\sqrt{\pm D}$. In this way only half of all the forms of the preceding article remain. We will denote their collection by $W$, and need only show how many distinct classes arise from these forms. Moreover, when $D$ is negative, there are clearly as many positive as negative forms in $W$.

I. When $D$ is negative, the individual forms in $W$ belong to distinct classes. For all the forms $(A,0,C)$ are reduced; likewise all the forms $(2B,B,C)$ are reduced except those in which $C<2B$. In such a form $2C<2B+C$, so (since $B<\frac{D}{B}$, \textit{that is,}\ $B<2C-B$, and therefore $2B<2C$, or $B<C$) we have $2C-2B<C$ and $C-B<\frac12 C$, and hence $(C,C-B,C)$, which is clearly equivalent to it, is a reduced form. This gives as many reduced forms as there are forms in $W$, and since it is easy to see that neither identical nor opposite forms can occur among them (with the single exception $C-B=0$, in which $B=C=\pm1$, and therefore $D=-1$, which we have already set aside), they all belong to distinct classes. We conclude that the number of all properly primitive ambiguous classes of determinant $D$ equals the number of forms in $W$, or half the number of forms in the preceding article. In the exceptional case $D=-1$ the same result holds by compensation: there are two classes, one containing the forms $(1,0,1)$, $(2,1,1)$, and the other $(-1,0,-1)$, $(-2,-1,-1)$. In general, then, for a negative determinant the number of all properly primitive ambiguous classes equals the number of all
\SourceLeaf{293}
assignable characters of primitive forms of this determinant; but the number of positive properly primitive ambiguous classes is half as great.

II. When $D$ is a positive square $=hh$, it is not difficult to prove that the individual forms in $W$ belong to distinct classes. But in this case we can reach the solution still more briefly as follows. By Article 210 every properly primitive ambiguous class of determinant $hh$, and no other class, contains one reduced form $(a,h,0)$, in which $a$ is a value of the expression $\sqrt{1}\ (\mathrm{mod.}\ 2h)$ lying between $0$ and $2h-1$ inclusive. Thus there are as many properly primitive ambiguous classes of determinant $hh$ as that expression has values. From Article 105 it follows without difficulty that the number of these values is $2^n$ or $2^{n+1}$ or $2^{n+2}$, according as $h$ is odd, twice an odd number, or divisible by four; that is, according as $D\equiv1$ or $\equiv4$ or $\equiv0\ (\mathrm{mod.}\ 8)$, where $n$ denotes the number of distinct odd prime divisors of $h$, or of $D$. Thus the number of properly primitive ambiguous classes is always half the total number of forms found in the preceding article, or equal to the number of forms in $W$, or of all possible characters.

III. When $D$ is a positive nonsquare, derive from each form $(A,B,C)$ in $W$ another form $(A,B',C')$ by taking $B'\equiv B\ (\mathrm{mod.}\ A)$ between the bounds $\sqrt{D}$ and $\sqrt{D}\mp A$ (using the upper or lower sign according as $A$. is positive or negative), and $C'=\frac{B'B'-D}{A}$; denote their collection by $W'$. Clearly these forms are properly primitive ambiguous forms of determinant $D$, all distinct; moreover, they are all reduced. For when $A<\sqrt{D}$, $B'$ is clearly $<\sqrt{D}$ and positive; also $B'>\sqrt{D}\mp A$, so $A>\sqrt{D}-B'$, and therefore $A$, taken positively, certainly lies between $\sqrt{D}+B'$ and $\sqrt{D}-B'$. But when $A>\sqrt{D}$, we cannot have $B=0$ (since we discarded those forms), and must have $B=\frac12 A$. Thus $B'$ equals $\frac12 A$ in magnitude and has a positive sign (for since $A<2\sqrt{D}$, $\pm\frac12 A$ lies between the bounds assigned to $B'$ and is congruent to $B$ modulo $A$, so $B'=\pm\frac12 A$). Hence $B'<\sqrt{D}$, which gives $2B'<\sqrt{D}+B'$, or $A<\sqrt{D}+B'$; therefore $\pm A$ must lie between the bounds $\sqrt{D}+B'$ and $\sqrt{D}-B'$. Finally, $W'$ contains all reduced properly primitive ambiguous forms of determinant $D$. For if $(a,b,c)$ is such a form, either $b\equiv0$ or $b\equiv\frac12 a\ (\mathrm{mod.}\ a)$. In the former case we clearly cannot have $b<a$, and therefore not $a>\sqrt{D}$, so the form $\left(a,0,-\frac{D}{a}\right)$ certainly belongs
\SourceLeaf{294}
to $W$, and the corresponding $(a,b,c)$ to $W'$; in the latter we certainly have $a<2\sqrt{D}$, so $\left(a,\allowbreak \frac12 a,\allowbreak \frac14 a-\frac{D}{a}\right)$ belongs to $W$, and the corresponding $(a,b,c)$ to $W'$. Thus the number of forms in $W$ equals the number of all reduced properly primitive ambiguous forms of determinant $D$. But since each ambiguous class contains \textit{two} reduced ambiguous forms (Articles 187, 194), the number of all properly primitive ambiguous classes of determinant $D$ is half the number of forms in $W$, or half the number of all assignable characters.

\Needspace{8\baselineskip}
\Article{259}
The number of improperly primitive ambiguous classes of a given determinant $D$ always equals the number of properly primitive ambiguous classes of the same determinant. Let $K$ be the principal class, and $K'$, $K''$ etc. the remaining properly primitive ambiguous classes of this determinant; let $L$ be any improperly primitive ambiguous class of the same determinant, \textit{for example,}\ the one containing the form $(2,1,\frac12-\frac12 D)$. Composing the class $L$ with $K$ gives the class $L$ itself. Suppose composing the class $L$ with $K'$, $K''$ etc. gives the classes $L'$, $L''$ etc. respectively. Clearly all belong to the same determinant $D$ and are improperly primitive and ambiguous. The theorem will therefore be proved once we have shown that all the classes $L$, $L'$, $L''$ etc. are distinct and that there are no other improperly primitive ambiguous classes of determinant $D$. For this we distinguish the following cases.

I. When the number of improperly primitive classes equals the number of properly primitive classes, each of the former arises by composing the class $L$ with a definite properly primitive class, so all $L$, $L'$, $L''$ etc. must be distinct. If $\mathfrak{L}$ denotes any improperly primitive ambiguous class of determinant $D$, there is a properly primitive class $\mathfrak{K}$ such that $\mathfrak{K}+L=\mathfrak{L}$. If the class $\mathfrak{K}$ has opposite class $\mathfrak{K}'$, we also have (since the classes $L$, $\mathfrak{L}$ are their own opposites) $\mathfrak{K}'+L=\mathfrak{L}$. Thus $\mathfrak{K}$ must be identical to $\mathfrak{K}'$ and therefore ambiguous; so $\mathfrak{K}$ occurs among the classes $K$, $K'$, $K''$ etc., and $\mathfrak{L}$ among $L$, $L'$, $L''$ etc.

II. When the number of improperly primitive classes is one third the number of properly primitive classes, let $H$ be the class containing the form $\left(4,\allowbreak 1,\allowbreak \frac{1-D}{4}\right)$, and $H'$ the one containing $\left(4,\allowbreak 3,\allowbreak \frac{9-D}{4}\right)$. Then $H$, $H'$ are properly primitive and distinct both
\SourceLeaf{295}
from one another and from the principal class $K$, and $H+H'=K$; $2H=H'$, $2H'=H$. If $\mathfrak{L}$ is any improperly primitive class of determinant $D$ arising by composing the class $L$ with the properly primitive class $\mathfrak{K}$, we also have $\mathfrak{L}=L+\mathfrak{K}+H$ and $\mathfrak{L}=L+\mathfrak{K}+H'$. Besides the three (properly primitive and distinct) classes $\mathfrak{K}$, $\mathfrak{K}+H$, $\mathfrak{K}+H'$, no others composed with $L$ produce $\mathfrak{L}$. Thus, if $\mathfrak{L}$ is ambiguous and $\mathfrak{K}'$ opposite to $\mathfrak{K}$, we also have $L+\mathfrak{K}'=\mathfrak{L}$, so $\mathfrak{K}'$ must be identical to one of those three classes. If $\mathfrak{K}'=\mathfrak{K}$, $\mathfrak{K}$ is ambiguous; if $\mathfrak{K}'=\mathfrak{K}+H$, we have $K=\mathfrak{K}+\mathfrak{K}'=2\mathfrak{K}+H=2(\mathfrak{K}+H')$, so $\mathfrak{K}+H'$ is ambiguous; likewise if $\mathfrak{K}'=\mathfrak{K}+H'$, $\mathfrak{K}+H$ is ambiguous. We conclude that $\mathfrak{L}$ must occur among the classes $L$, $L'$, $L''$ etc. But it is easy to see that among the three classes $\mathfrak{K}$, $\mathfrak{K}+H$, $\mathfrak{K}+H'$ more than one cannot be ambiguous. For if both $\mathfrak{K}$ and $\mathfrak{K}+H$ were ambiguous, or identical to their respective opposites $\mathfrak{K}'$, $\mathfrak{K}'+H'$, we would have $\mathfrak{K}+H=\mathfrak{K}+H'$. The same conclusion follows from assuming $\mathfrak{K}$ and $\mathfrak{K}+H'$ ambiguous; finally, if $\mathfrak{K}+H$, $\mathfrak{K}+H'$ were ambiguous, or identical to their opposites $\mathfrak{K}'+H'$, $\mathfrak{K}'+H$, we would have $\mathfrak{K}+H+\mathfrak{K}'+H=\mathfrak{K}'+H'+\mathfrak{K}+H'$, giving $2H=2H'$, or $H'=H$. Thus there is only one properly primitive ambiguous class whose composition with $L$ produces $\mathfrak{L}$, so all $L$, $L'$, $L''$ etc. are distinct.

The number of ambiguous classes in a \textit{derived} order clearly equals the number of ambiguous classes in the primitive order from which it is derived, and so can always be determined from what precedes.

\Needspace{8\baselineskip}
\Article{260}
\textsc{Problem.} \textit{A properly primitive class $K$ of determinant $D$ arises by doubling a properly primitive class $k$ of the same determinant. Find a\,l\,l such classes whose doubling gives the class $K$.}

\textit{Solution.} Let $H$ be the principal class of determinant $D$, and $H'$, $H''$, $H'''$ etc. the remaining properly primitive ambiguous classes of the same determinant; denote the classes obtained by composing these with $k$, namely $k+H'$, $k+H''$, $k+H'''$, by $k'$, $k''$, $k'''$ etc. Then all the classes $k$, $k'$, $k''$ etc. are properly primitive, of determinant $D$, and distinct; it is just as easy to see that doubling each gives the class $K$. If $\mathfrak{K}$ denotes any properly primitive class of determinant $D$ whose doubling produces the class $K$, it must occur among the classes $k$, $k'$, $k''$ etc. For set $\mathfrak{K}=k+\mathfrak{H}$, where $\mathfrak{H}$ is a properly primitive class of determinant $D$ (Article 249). and then
\SourceLeaf{296}
$2k+2\mathfrak{H}=2\mathfrak{K}=K=2k$, from which we easily conclude that $2\mathfrak{H}$ coincides with the principal class, $\mathfrak{H}$ is ambiguous, that is, occurs among $H$, $H'$, $H''$ etc., and $\mathfrak{K}$ among $k$, $k'$, $k''$ etc. These classes therefore give the complete solution to the problem.

Moreover, when $D$ is negative, clearly half of the classes $k$, $k'$, $k''$ etc. are positive and half negative.

Thus every properly primitive class of determinant $D$ that can arise by doubling any such class arises by doubling exactly as many such classes as there are properly primitive ambiguous classes of determinant $D$. Consequently, if the number of all properly primitive classes of determinant $D$ is $r$, and the number of all properly primitive ambiguous classes of this determinant is $n$, the number of all properly primitive classes of the same determinant obtainable by doubling such a class is $\frac{r}{n}$. The same formula results if, for a negative determinant, the symbols $r$, $n$ denote the numbers of \textit{positive} classes: the former \textit{all} properly primitive classes, the latter only the ambiguous ones. Thus, \textit{for example,}\ for $D=-161$ the number of all positive properly primitive classes is $16$, and the number of ambiguous classes $4$, so the number of all classes obtainable by doubling some class must be $4$. And indeed all classes in the principal genus are found to have this property: the principal class $(1,0,161)$ arises by doubling the four ambiguous classes; $(2,1,81)$ by doubling the classes $(9,1,18)$, $(9,-1,18)$, $(11,2,15)$, $(11,-2,15)$; $(9,1,18)$ by doubling the classes $(3,1,54)$, $(6,1,27)$, $(5,-2,33)$, $(10,3,17)$; finally $(9,-1,18)$ by doubling the classes $(3,-1,54)$, $(6,-1,27)$, $(5,2,33)$, $(10,-3,17)$.
